\section{Background and Related Work}
\label{sec:background}

\paragraph{Alignment and refusal.}
Instruction tuning and reinforcement learning from human or AI feedback make
language models more helpful and less likely to produce harmful content
\cite{christiano2017deep,ziegler2019fine,stiennon2020learning,
askell2021general,ouyang2022training,bai2022constitutional}. In deployed
assistants, this often appears as refusal behavior: the model declines to
complete requests it classifies as unsafe. Refusal is a useful safety mechanism,
but it is a context-poor one when judged only from the prompt text. Security
work creates many cases where the natural-language description resembles an
attack even though the environment is authorized, local, and intentionally
vulnerable. Prior safety evaluations show that this helpfulness-harmlessness
tension can produce both under-refusal and over-refusal
\cite{lin2021truthfulqa,hendrycks2021ethics,weidinger2021ethical,
rotter2023xstest} in deployed systems.

\paragraph{Robust refusal and jailbreak benchmarks.}
HarmBench and JailbreakBench standardize the evaluation of adversarial prompts,
harmful behaviors, and robust refusal \cite{mazeika2024harmbench,
chao2024jailbreakbench}. StrongREJECT, XSTest, and universal prompt-attack
work further clarify how refusal scoring, false refusal, and jailbreak
transferability should be measured \cite{souly2024strongreject,
rotter2023xstest,zou2023universal,wei2023jailbroken}. These benchmarks are
valuable because they clarify threat models, scoring rules, and attack
artifacts. Our evaluation differs in the target behavior. We do not ask whether
a jailbreak can make a model provide harmful instructions in open-ended chat. We
ask how refusal-oriented alignment changes task completion and evidence quality
when an agent must act over many turns inside a controlled security environment.

\paragraph{Cybersecurity capability benchmarks.}
CyberSecEval 2, CyberSecEval 3, Cybench, NYUCTF Bench, EnIGMA, and SEC-bench
evaluate cybersecurity risks and capabilities of language models, including
prompt injection, exploit generation, vulnerability patching, and
capture-the-flag style tasks \cite{bhatt2024cyberseceval,li2024cyberseceval3,
zhang2025cybench,shao2024nyuctf,abramovich2024enigma,
secbench2025automated}. These efforts show the importance of executable
environments, task-specific success predicates, and partial-credit
decomposition. We build on that benchmarking style but center paired
comparisons that keep the harness, task, and seed fixed within each evaluated
pair while varying the alignment condition.

\paragraph{Agent benchmarks and tool use.}
AgentBench, ReAct, Toolformer, WebShop, WebArena, InterCode, SWE-bench,
SWE-agent, and OSWorld show that multi-turn tool use introduces failure modes
that are invisible in single-turn question answering
\cite{li2023agentbench,yao2023react,schick2023toolformer,yao2022webshop,
zhou2023webarena,yang2023intercode,jimenez2024swebench,yang2024sweagent,
xie2024osworld}. Recent critiques emphasize that agent benchmarks require
careful success predicates, environment controls, and resistance to benchmark
gaming \cite{kapoor2024aiagents}. These lessons motivate our use of
machine-readable task manifests, deterministic checkers, and complete traces.

\looseness=-1
\paragraph{Defensive refusal in cybersecurity.}
Recent work on \emph{defensive refusal bias} \cite{campbell2026defensive}
studies the failure mode closest to ours. On 2,390 single-turn conversations
from the National Collegiate Cyber Defense Competition, Claude 3.5 Sonnet,
GPT-4o, and Llama-3.3-70B-Instruct together refuse defensive requests containing
security-sensitive keywords 2.72 times as often as semantically equivalent
neutral requests, and explicit authorization is associated with more refusals,
not fewer. Where that study classifies refusals in response text, we evaluate
multi-turn autonomous execution in executable sandboxes with deterministic
success predicates, saved traces, failure phases, unsafe-action labels, and
evidence grounding, comparing aligned models with their less-restricted
derivatives. Both studies find that authorization wording is not a
reliable lever: our authorized Gemma 4 31B condition succeeds in 0/150 runs.
But none of our security traces contains a refusal event; the aligned Gemma
models keep working and produce less-grounded artifacts, a difference that
refusal rates miss.

\paragraph{Mechanisms of refusal removal.}
Recent mechanistic work suggests that refusal behavior in open-weight chat
models can sometimes be controlled through low-dimensional activation
directions \cite{arditi2024refusal}. Community ``uncensored'' or
``abliterated'' models apply related ideas to reduce refusals. Such derivatives
are useful for paired comparisons, but they are not perfect counterfactuals: weight
editing, quantization, template changes, or fine-tuning may alter more than
refusal behavior. Our protocol therefore requires matched inference settings,
non-security sanity checks, and explicit limitations around derivative-model
provenance.

\paragraph{Model families and derivative provenance.}
Gemma 4 is the central family because the official model cards describe
instruction-tuned open-weight models with long context windows, coding ability,
structured-output support, and agentic use cases
\cite{google_gemma4_blog_2026,hf_gemma4_31b_it,hf_gemma4_26b_a4b_it}. The
TrevorJS derivatives advertise reduced refusal behavior for the same Gemma 4
family \cite{trevorjs_gemma4_uncensored_collection,
trevorjs_gemma4_31b_uncensored,trevorjs_gemma4_26b_uncensored}. The 26B
evaluation uses local GGUF quantizations from Unsloth and TrevorJS
\cite{unsloth_gemma4_26b_a4b_gguf,trevorjs_gemma4_26b_uncensored_gguf}.

The cross-family evaluation also includes Qwen2.5-Coder 7B and Llama 3.1 8B.
For each additional family, the aligned condition uses an instruction-tuned
GGUF release and the comparison condition uses an abliterated derivative
distributed as a GGUF quantization
\cite{qwen25coder7b_gguf,qwen25coder7b_abliterated_gguf,
llama31_8b_instruct_gguf,llama31_8b_abliterated_gguf}. These are not clean
causal counterfactuals for safety alignment. We therefore report them as
cross-family comparisons and treat model provenance as part of the
interpretation rather than as a nuisance detail.

\looseness=-1
The derivatives also remove refusals in different ways, so we attribute
effects to these specific transformations rather than to abliteration in
general. According to the model cards, the Gemma 4 31B derivative uses
norm-preserving biprojected abliteration \cite{grimjim2025normpreserving},
projecting per-layer refusal directions out of the attention-output and MLP
down-projection weights of every layer while preserving row norms
\cite{trevorjs_gemma4_31b_uncensored}; the 26B A4B derivative additionally
applies expert-granular abliteration to all 128 expert down-projections in each
mixture-of-experts layer \cite{trevorjs_gemma4_26b_uncensored}. The Llama 3.1
derivative uses refusal-direction abliteration \cite{arditi2024refusal} via
FailSpy's implementation \cite{mlabonne_llama31_8b_abliterated}; the
Qwen2.5-Coder derivative is an abliterated release by huihui-ai whose GGUF card
does not document the procedure \cite{qwen25coder7b_abliterated_gguf}. None
isolates safety fine-tuning: public base checkpoints precede instruction tuning
altogether, and no official instruction-tuned checkpoints of these models
without safety training are available.
